% Generated by scripts/assemble.py; corrections belong in source/corrections.json.
\examhead{85}{哈尔滨工程大学 2025 年数学分析试卷(回忆版)}{84}
% p085-b01; source PDF page 85
\blockanchor{p085-b01}
\question{1}计算 $\lim_{x\to 0}\left(\frac{1}{x^2} -\frac{1}{(\tan x)^2}\right)$.\par
% p085-b02; source PDF page 85
\blockanchor{p085-b02}
\question{2}证明: 对任意的 $x > 0$, 有 $0 < \frac{1}{\ln(1 + x)} - \frac{1}{x} < 1$.\par
% p085-b03; source PDF page 85
\blockanchor{p085-b03}
\question{3}讨论下列级数的敛散性 $\sum_{n=3}^{\infty} \frac{1}{n(\ln n)^p (\ln \ln n)^q}$, 其中 $p, q$ 为实数.\par
% p085-b04; source PDF page 85
\blockanchor{p085-b04}
\question{4}设 $f(x)$ 在 $x = 0$ 处连续, 且对任意的 $x, y \in \mathbb{R}$, 有 $f(x + y) = f(x) + f(y)$. 证明:\par
% p085-b05; source PDF page 85
\blockanchor{p085-b05}
(1) $f$ 在 $\mathbb{R}$ 上连续.\par
% p085-b06; source PDF page 85
\blockanchor{p085-b06}
(2) 存在实数 $M$, 使得 $|f(x)| \leq M|x|$.\par
% p085-b07; source PDF page 85
\blockanchor{p085-b07}
\question{5}求积分 $\int_0^1\frac{\ln x}{\sqrt{1 - x^2}}\mathrm{d}x.$\par
% p085-b08; source PDF page 85
\blockanchor{p085-b08}
\question{6}设函数 $u(x, y) = \sum_{n=1}^{\infty} \sum_{m=1}^{\infty} A_{mn} \sin(m\pi x) \sin(n\pi y)$ 满足 $u_{xx} + u_{yy} = x, (x, y) \in [0, 1] \times [0, 1]$, 求 $A_{mn}$.\par
% p085-b09; source PDF page 85
\par\Needspace{6\baselineskip}
\blockanchor{p085-b09}
\question{7}令函数 $f$ 为定义在 $\mathbb{R}$ 上的连续函数, 证明: 对任意的 $x, y, z \in \mathbb{R}^{+}$, 下式成立:\par
\nopagebreak[4]
% p085-b10; source PDF page 85
\blockanchor{p085-b10}
\displayformula{\int_ {0} ^ {x} \int_ {0} ^ {y} \int_ {0} ^ {z} f (t) \mathrm{d} t \mathrm{d} z \mathrm{d} y = \frac {1}{2} \int_ {0} ^ {x} (x - t) ^ {2} f (t) \mathrm{d} t.}
% p085-b11; source PDF page 85
\blockanchor{p085-b11}
\question{8}求函数 $f(x, y) = (x^2 + 2y^2)e^{-(x^2 + y^2)}$ 在区域 $D = \{(x, y) | x^2 + y^2 \leq 4\}$ 上的最大值和最小值.\par
% p085-b12; source PDF page 85
\blockanchor{p085-b12}
\question{9}计算二重积分 $\iint_{R} (x^{2} + y^{2}) \, \mathrm{d}x \, \mathrm{d}y$, 其中 $R = \{(x, y) \mid x^{2} - xy + y^{2} \leq 2\}$.\par
% p085-b13; source PDF page 85
\par\Needspace{6\baselineskip}
\blockanchor{p085-b13}
\question{10}设 $B = B(\mathbf{0}, 1)$ 是 $\mathbb{R}^3$ 中以原点为中心以 1 为半径的开球体, $\partial B$ 为其表面. 令\par
\nopagebreak[4]
% p085-b14; source PDF page 85
\blockanchor{p085-b14}
\displayformula{\boldsymbol {u} = (u (x, y, z), v (x, y, z), w (x, y, z))}
% p085-b15; source PDF page 85
\par\Needspace{6\baselineskip}
\blockanchor{p085-b15}
为 $\mathbb{R}^3$ 中的向量场, 满足 $u_x + v_y + w_z = 0$. 且在 $\partial B$ 中满足 $\pmb{u} \cdot \pmb{n} = 0$, 其中 $\pmb{n}$ 为 $\partial B$ 的单位外法向量. 证明下式成立:\par
\nopagebreak[4]
% p085-b16; source PDF page 85
\blockanchor{p085-b16}
\displayformula{\iiint_ {B} [ (\boldsymbol {u} \cdot \nabla) \boldsymbol {u} ] \cdot \boldsymbol {u} d x d y d z = 0.}
% p085-b17; source PDF page 85
\blockanchor{p085-b17}
其中  $u \cdot \nabla = u \frac{\partial}{\partial x} + v \frac{\partial}{\partial y} + w \frac{\partial}{\partial z}.$\par
